\section{Norms on finite-dimensional vector spaces}

\begin{definition}[Norm over a multiplicatively valued field]
    Let \((K,\varphi)\) be a field equipped with a multiplicative valuation, and let
    \(V\) be a vector space over \(K\). A function
    \[
        \|\cdot\|:V\longrightarrow \mathbb{R}_{\geq 0}
    \]
    is called a non-archimedean norm on \(V\) compatible with \(\varphi\) if the
    following conditions hold:
    \begin{enumerate}
        \item For every \(v\in V\),
              \[
                  \|v\|\geq 0,
              \]
              and
              \[
                  \|v\|=0 \Longleftrightarrow v=0.
              \]

        \item For every \(\alpha\in K\) and every \(v\in V\),
              \[
                  \|\alpha v\|=\varphi(\alpha)\|v\|.
              \]

        \item For every \(u,v\in V\),
              \[
                  \|u+v\|\leq \max\{\|u\|,\|v\|\}.
              \]
    \end{enumerate}
\end{definition}

\begin{proposition}
    Let \(V\) be a vector space over a multiplicatively valued field \((K,\varphi)\).
    Suppose that \(V\) is equipped with a non-archimedean norm \(\|\cdot\|\).
    Then \(V\) becomes a metric space with metric
    \[
        d(u,v):=\|u-v\|.
    \]
    Moreover, the maps
    \[
        V\times V\longrightarrow V,
        \qquad
        (u,v)\longmapsto u+v,
    \]
    and
    \[
        K\times V\longrightarrow V,
        \qquad
        (\alpha,v)\longmapsto \alpha v,
    \]
    are both continuous.
\end{proposition}

\begin{proof}
    The function \(d(u,v)=\|u-v\|\) is a metric. Indeed, for \(u,v\in V\),
    \[
        d(u,v)=0
        \Longleftrightarrow
        \|u-v\|=0
        \Longleftrightarrow
        u=v.
    \]
    Also,
    \[
        d(u,v)=\|u-v\|=\|-(v-u)\|=\|v-u\|=d(v,u),
    \]
    because \(\varphi(-1)=1\). Finally, for \(u,v,w\in V\),
    \[
        d(u,w)
        =
        \|u-w\|
        =
        \|(u-v)+(v-w)\|
        \leq
        \max\{\|u-v\|,\|v-w\|\}
        =
        \max\{d(u,v),d(v,w)\}.
    \]
    Hence \(d\) is a non-archimedean metric on \(V\).

    The addition map is continuous because
    \[
        \|(u+v)-(u_0+v_0)\|
        =
        \|(u-u_0)+(v-v_0)\|
        \leq
        \max\{\|u-u_0\|,\|v-v_0\|\}.
    \]

    The scalar multiplication map is continuous because
    \[
        \|\alpha v-\alpha_0v_0\|
        =
        \|\alpha(v-v_0)+(\alpha-\alpha_0)v_0\|
        \leq
        \max\{\varphi(\alpha)\|v-v_0\|,
               \varphi(\alpha-\alpha_0)\|v_0\|\}.
    \]
    Since \(\varphi(\alpha)\) is locally bounded near \(\alpha_0\), the right-hand side
    tends to \(0\) whenever \(\alpha\to\alpha_0\) in \(K\) and \(v\to v_0\) in \(V\).
    Thus scalar multiplication is continuous.
\end{proof}

\begin{lemma}
    Let \((K,\varphi)\) be a multiplicatively valued field, and let \(V\) be a
    finite-dimensional \(K\)-vector space with \(K\)-basis
    \[
        u_1,\dots,u_n.
    \]
    For
    \[
        v=a_1u_1+\cdots+a_nu_n\in V,
    \]
    define
    \[
        \|v\|_0:=\max\{\,\varphi(a_i)\mid 1\leq i\leq n\,\}.
    \]
    Then \(\|\cdot\|_0\) is a non-archimedean norm on \(V\). Moreover, the following
    hold:
    \begin{enumerate}
        \item Let
              \[
                  v_r=a_{r,1}u_1+\cdots+a_{r,n}u_n\in V
                  \qquad (r\geq 1).
              \]
              Then the sequence \((v_r)\) is a Cauchy sequence in \(V\) if and only if
              each coordinate sequence \((a_{r,i})_{r\geq 1}\) is a Cauchy sequence in
              \(K\).

              Also, for
              \[
                  v=a_1u_1+\cdots+a_nu_n\in V,
              \]
              one has
              \[
                  \lim_{r\to\infty} v_r=v
                  \Longleftrightarrow
                  \lim_{r\to\infty} a_{r,i}=a_i
                  \quad\text{for all }1\leq i\leq n.
              \]

        \item If \(K\) is complete, then \(V\) is complete with respect to the norm
              \(\|\cdot\|_0\).
    \end{enumerate}
\end{lemma}

\begin{proof}
    First we verify that \(\|\cdot\|_0\) is a norm. Clearly
    \[
        \|v\|_0\geq 0.
    \]
    Moreover,
    \[
        \|v\|_0=0
        \Longleftrightarrow
        \varphi(a_i)=0\quad\text{for all }i
        \Longleftrightarrow
        a_i=0\quad\text{for all }i
        \Longleftrightarrow
        v=0.
    \]
    For \(\lambda\in K\), we have
    \[
        \|\lambda v\|_0
        =
        \max_i \varphi(\lambda a_i)
        =
        \varphi(\lambda)\max_i\varphi(a_i)
        =
        \varphi(\lambda)\|v\|_0.
    \]
    Finally, if
    \[
        v=a_1u_1+\cdots+a_nu_n,\qquad
        w=b_1u_1+\cdots+b_nu_n,
    \]
    then
    \[
        v+w=(a_1+b_1)u_1+\cdots+(a_n+b_n)u_n.
    \]
    Hence
    \[
        \|v+w\|_0
        =
        \max_i\varphi(a_i+b_i)
        \leq
        \max_i\max\{\varphi(a_i),\varphi(b_i)\}
        =
        \max\{\|v\|_0,\|w\|_0\}.
    \]
    Thus \(\|\cdot\|_0\) is a non-archimedean norm.

    Now let
    \[
        v_r=a_{r,1}u_1+\cdots+a_{r,n}u_n.
    \]
    For \(r,s\geq 1\), we have
    \[
        v_r-v_s
        =
        \sum_{i=1}^n (a_{r,i}-a_{s,i})u_i.
    \]
    Therefore
    \[
        \|v_r-v_s\|_0
        =
        \max_{1\leq i\leq n}
        \varphi(a_{r,i}-a_{s,i}).
    \]
    It follows immediately that \((v_r)\) is Cauchy in \(V\) if and only if each
    \((a_{r,i})_{r\geq 1}\) is Cauchy in \(K\).

    Similarly,
    \[
        v_r-v
        =
        \sum_{i=1}^n (a_{r,i}-a_i)u_i,
    \]
    so
    \[
        \|v_r-v\|_0
        =
        \max_{1\leq i\leq n}
        \varphi(a_{r,i}-a_i).
    \]
    Hence
    \[
        \lim_{r\to\infty} v_r=v
        \Longleftrightarrow
        \lim_{r\to\infty} a_{r,i}=a_i
        \quad\text{for all }i.
    \]

    Finally, assume that \(K\) is complete. Let \((v_r)\) be a Cauchy sequence in
    \(V\). By the first part, each coordinate sequence \((a_{r,i})\) is Cauchy in
    \(K\). Since \(K\) is complete, there exists \(a_i\in K\) such that
    \[
        \lim_{r\to\infty}a_{r,i}=a_i.
    \]
    Put
    \[
        v=a_1u_1+\cdots+a_nu_n.
    \]
    By the convergence criterion above,
    \[
        \lim_{r\to\infty}v_r=v.
    \]
    Therefore \(V\) is complete.
\end{proof}

\begin{theorem}[Equivalence of norms over a complete non-archimedean field]
    Let \((K,\varphi)\) be a complete non-archimedean normed field, and let
    \(V\) be a finite-dimensional \(K\)-vector space with basis
    \[
        u_1,\dots,u_n.
    \]
    Define the associated supremum norm by
    \[
        \left\|\sum_{i=1}^{n} a_i u_i\right\|_{\infty}
        :=
        \max_{1\leq i\leq n}\varphi(a_i).
    \]
    Then, for any non-archimedean norm \(\|\cdot\|\) on \(V\), there exist constants
    \(\lambda,\mu>0\) such that
    \[
        \|v\|\leq \lambda \|v\|_{\infty},
        \qquad
        \|v\|_{\infty}\leq \mu \|v\|
    \]
    for all \(v\in V\). Consequently, \(\|\cdot\|\) and \(\|\cdot\|_{\infty}\) induce the same
    topology on \(V\). In particular, any two norms on \(V\) induce the same topology.
\end{theorem}

\begin{proof}
    Let
    \[
        \lambda=\max_{1\leq i\leq n}\|u_i\|.
    \]
    For
    \[
        v=\sum_{i=1}^{n}a_i u_i,
    \]
    the non-archimedean triangle inequality gives
    \[
        \|v\|
        \leq
        \max_{1\leq i\leq n}\|a_i u_i\|
        =
        \max_{1\leq i\leq n}\varphi(a_i)\|u_i\|
        \leq
        \lambda \|v\|_{\infty}.
    \]
    Hence the first inequality holds.

    It remains to prove that there exists \(\mu>0\) such that
    \[
        \|v\|_{\infty}\leq \mu \|v\|
    \]
    for all \(v\in V\). We prove this by induction on \(n=\dim_K V\).

    If \(n=1\), then every \(v\in V\) can be written as \(v=a_1u_1\). Hence
    \[
        \|v\|=\varphi(a_1)\|u_1\|,
        \qquad
        \|v\|_{\infty}=\varphi(a_1).
    \]
    Therefore
    \[
        \|v\|_{\infty}
        =
        \frac{1}{\|u_1\|}\|v\|,
    \]
    so the claim holds.

    Now assume \(n>1\), and suppose the result is known for all vector spaces of
    dimension \(n-1\). For
    \[
        v=\sum_{i=1}^{n}a_i u_i,
    \]
    define
    \[
        \psi_i(v):=\varphi(a_i).
    \]
    We claim that, for each \(i\), there exists a constant \(\mu_i>0\) such that
    \[
        \psi_i(v)\leq \mu_i\|v\|
    \]
    for all \(v\in V\). Once this is proved, taking
    \[
        \mu=\max_{1\leq i\leq n}\mu_i
    \]
    gives
    \[
        \|v\|_{\infty}
        =
        \max_{1\leq i\leq n}\psi_i(v)
        \leq
        \mu\|v\|.
    \]

    It suffices to prove the claim for \(i=n\), since the other cases follow by
    reordering the basis. Put
    \[
        V' = Ku_1\oplus\cdots\oplus Ku_{n-1}.
    \]
    By the induction hypothesis, the restrictions of \(\|\cdot\|\) and
    \(\|\cdot\|_{\infty}\) to \(V'\) induce the same topology.

    Since \(K\) is complete, \(V'\) is complete with respect to the norm
    \(\|\cdot\|_{\infty}\). Hence \(V'\) is also complete with respect to
    \(\|\cdot\|\). Therefore \(V'\) is closed in \(V\) with respect to the topology
    induced by \(\|\cdot\|\).

    Suppose, for contradiction, that there is no constant \(\mu_n>0\) such that
    \[
        \psi_n(v)\leq \mu_n\|v\|
    \]
    for all \(v\in V\). Then, for every integer \(j\geq 1\), there exists
    \(v_j\in V\) such that
    \[
        \psi_n(v_j)>j\|v_j\|.
    \]
    Since the inequality is homogeneous, we may multiply \(v_j\) by a nonzero scalar
    and assume that the coefficient of \(u_n\) in \(v_j\) is \(1\). Thus we may write
    \[
        v_j=a_{1j}u_1+\cdots+a_{n-1,j}u_{n-1}+u_n.
    \]
    Then
    \[
        \psi_n(v_j)=1,
    \]
    and the inequality above gives
    \[
        \|v_j\|<\frac{1}{j}.
    \]
    Hence
    \[
        v_j\longrightarrow 0
    \]
    with respect to \(\|\cdot\|\).

    On the other hand,
    \[
        v_j-u_n\in V'
    \]
    for every \(j\), and
    \[
        v_j-u_n\longrightarrow -u_n
    \]
    with respect to \(\|\cdot\|\). Since \(V'\) is closed in \(V\), it follows that
    \[
        -u_n\in V'.
    \]
    This is impossible, because \(u_1,\dots,u_n\) are linearly independent.

    Therefore there exists \(\mu_n>0\) such that
    \[
        \psi_n(v)\leq \mu_n\|v\|
    \]
    for all \(v\in V\). As explained above, the same argument applies to each
    coordinate \(\psi_i\). Hence there exists \(\mu>0\) such that
    \[
        \|v\|_{\infty}\leq \mu\|v\|
    \]
    for all \(v\in V\).

    Combining the two inequalities, we conclude that \(\|\cdot\|\) and
    \(\|\cdot\|_{\infty}\) are equivalent norms. Therefore they induce the same
    topology on \(V\).
\end{proof}
